\section{A Relational Calculus in ZFC}
\label{sec:zflean-rel}

Throughout this section we work inside the ZFC model provided by Mathlib and extend it with a calculus for binary relations and partial/total functions.
A \emph{binary relation} $R$ between $A$ and $B$ is a subset $R \subseteq A\times B$.
Our calculus packages standard operations on relations (converse, identity, composition, domain, range, image) together with predicates for partial and total functions, and supplies small \emph{automation tactics} that try to discharge ubiquitous well-formedness side-conditions, so that these objects can be manipulated \emph{as is} without cumbersome boilerplate, as on paper.
Those tactics are called \inlinelean|zrel|, \inlinelean|zpfun|, and \inlinelean|zfun|, and will be described later.

\subsection{Basic definitions}

\begin{definition}[converse, identity, composition]\label{def:rel-basic}
  Let $R\subseteq A\times B$ and $S\subseteq B\times C$ be relations.
  \begin{itemize}
    \item The \emph{identity relation} $\Id{A}$ on $A$ is defined as:
          \begin{equation*}
            \Id{A} \triangleq \{ x \in A \times A \mid \pi_1(x) = \pi_2(x) \}
          \end{equation*}
          The identity relation is shown to be a total function belonging to $A^A$, and later shown to be a bijection and the neutral element for composition.
    \item The \emph{converse} $R^{-1}$ is defined as:
          \begin{equation*}
            R^{-1} \triangleq \{z \in B\times A \mid \exists\, x\, y,\, x \in A \land y \in B \land z = (y, x) \land (x,y)\in R \}
          \end{equation*}
          This definition depends on the side-condition $R\subseteq A\times B$ (declaring that $R$ is a relation between the implicit sets), an implicit argument that \inlinelean|zrel| attempts to discharge in \zflean.
    \item The \emph{composition} $S\circ R \subseteq A\times C$ is defined as:
          \begin{equation*}
            S \circ R := \{ w \in A\times C \mid \exists\, x\, z,\, x \in A \land z \in C \land w = (x, z) \land \exists\, y \in B,\, (x,y)\in R \land (y, z)\in S \}
          \end{equation*}
          Unlike the previous definition, nothing is enforced on $R$ and $S$: if either $R$ or $S$ is not a relation, the composition is still well-defined, but is empty.
          However, we define \emph{functional composition} $\zbinop{\circ}$ specifically for functions, requiring $R$ and $S$ to be functions, and using the \inlinelean|zfun| tactic to try to discharge the side-conditions automatically.
  \end{itemize}
\end{definition}

\begin{definition}[Domain, range, image]\label{def:dom-range-image}
  For a relation $R\subseteq A\times B$, its \emph{domain}, \emph{range}, and \emph{image} of a set $X\subseteq A$ are defined as:
  \begin{align*}
    \begin{array}{rcl}
      \Dom(R)      & \triangleq & \{ x\in A \mid \exists\, y\in B,\, (x,y)\in R \}          \\
      \Ran(R)      & \triangleq & \{ y \in B \mid \exists\, x \in \Dom(R),\, (x,y) \in R \} \\
      \Image{R}{X} & \triangleq & \{ y \in B \mid \exists\, x \in X,\, (x, y) \in R \}
    \end{array}
  \end{align*}

  Operationally, all three definitions above invoke the \inlinelean|zrel| tactic to discharge the side-condition that $R$ is a relation between $A$ and $B$.
\end{definition}

Finally, we define partial functions as functional relations without totality requirements.
\begin{definition}[Partial functions]
  A set $f$ is said to be a partial function from set $A$ to set $B$, written $\mathsf{IsPFunc}(f,A,B)$, if it is \emph{functional}:\footnote{The order of arguments is voluntarily changed compared to $\IsFunc$, so that in Lean, one may write \inlinelean|f.IsPFunc A B|}
  \begin{equation*}
    \IsPFunc(f, A, B) \triangleq f \subseteq A \times B \,\land\, \left(\forall\, x \in A,\, \forall\, y \in B,\, \forall\, z \in B,\, (x,y) \in f \land (x,z) \in f \implies y = z\right)
  \end{equation*}
\end{definition}

We also derive predicates $\IsInjective$, $\IsSurjective$, and $\IsBijective$ in the usual sense, which all invoke the tactic \inlinelean|zfun| to discharge functional totality requirements.

\subsection{Function evaluation and λ-abstraction}

We introduce two convenient features for working with functions, namely function application and $\lambda$-abstraction, defined as follows.
We strive both to keep standard mathematical notation and to automate side-conditions as much as possible.
\begin{itemize}
  \item \emph{function application} that maps $x\in\Dom(f)$ to the unique $y$ such that $(x,y)\in f$ and denoted by $\fapply{f}{x}$, invoking the \inlinelean|zpfun| tactic; we show that evaluation commutes with composition in the expected way: $\fapply{(g \zbinop{\circ} f)}{x}=\fapply{g}{\fapply{f}{x}}$, and $\fapply{f}{x}$ coincides with the unique element of the image $f[\{x\}]$.
        In practice, $\fapplysym$ is a non-computable dependently typed operator that requires a proof that $f$ is a partial function and a proof $h_x$ that $x$ belongs to its domain (that must be provided, there is no automation at this point), denoted $\fapply{f}{\langle x, h_x \rangle}$.
        Its definition makes use of the axiom of choice to select an element $y$ such that $(x,y)\in f$.
        Since $f$ is functional, there is at most one such $y$, and since $x$ belongs to the domain of $f$, there is at least one such $y$, hence such a $y$ exists and is unique---this property is proved in \zflean{} as \leancodeptr{IsPFunc.exists\_unique\_of\_mem\_dom}{ZFLean/Functions.html\#ZFSet.IsPFunc.exists_unique_of_mem_dom}.

  \item $\lambda$-abstraction constructing functions from expressions, written as:
        \begin{equation*}
          \zlambda{A}{B}{x}{e(x)}
          \qquad\text{is the relation}\qquad
          \{ (x,y)\in A\times B \mid y=e(x) \},
        \end{equation*}
\end{itemize}
which is then shown to belong to $B^A$ whenever $e(x)\in B$ for all $x\in A$.
We prove the expected equational properties for this notation, summarized in the following theorem.

\begin{theorem}[Specification, extensionality, $\beta$- and $\eta$-conversion]
  \label{thm:lambda}
  Let $A$, $B$ be sets, $f \in B^A$, and $e,e'$ be unary Lean-level (endo)functions on sets.
  \begin{itemize}
    \item \textbf{Specification:} For any sets $x$ and $y$, the following holds:
          \begin{equation*}
            (x, y) \in \zlambda{A}{B}{x}{e(x)} \iff x \in A \land y \in B \land y = e(x)
          \end{equation*}
    \item \textbf{Extensionality:}
          \begin{equation*}
            \left(\zlambda{A}{B}{x}{e(x)}\right) = \left(\zlambda{A}{B}{x}{e'(x)}\right) \iff \forall\, x \in A,\, e(x) = e'(x)
          \end{equation*}
    \item \textbf{$\beta$-reduction:}
          \begin{equation*}
            \forall\, x \in A, \fapply{\left(\zlambda{A}{B}{x}{e(x)}\right)}{\langle x, \irrelprf \rangle} = e(x)
          \end{equation*}
    \item \textbf{$\eta$-expansion:}
          \begin{equation*}
            f = \left(\zlambda{A}{B}{x}{\fapply{f}{\langle x, \irrelprf \rangle}}\right)
          \end{equation*}
  \end{itemize}
\end{theorem}
\begin{proof}
  Proofs are carried out by applying extensionality on the corresponding sets.
  Full proof details are available in the artifacts as proofs of the theorems \inlinelean|lambda_spec|, \inlinelean|lambda_ext_iff|, \inlinelean|fapply_lambda|, and \inlinelean|lambda_eta|.
\end{proof}

Application and abstraction integrate well with rewriting, so that many proofs reduce to algebra on images, evaluation, and composition, as expected.
We illustrate this with a simple example below.

\begin{example}
  \label{ex:func-ext}
  Consider the extensionality principle for total functions: two functions $f,g\in B^A$ are equal if they are pointwise equal on $A$.
  In \zflean, this is stated as the following theorem:
  \begin{leanlisting}
theorem is_func_ext_iff {A B : ZFSet} {f g : ZFSet}
  (hf : IsFunc A B f) (hg : IsFunc A B g) :
  f = g ↔ ∀ x ∈ A, @ᶻf ⟨x, $\irrelprf$⟩ = @ᶻg ⟨x, $\irrelprf$⟩ := $\irrelprf$
  \end{leanlisting}
  where the first two omitted proofs ($\irrelprf$) are proofs that $x$ belongs to the domain of $f$ and $g$ respectively.
  The proof of this theorem can be carried out directly using extensionality on sets, although it turns out to be rather tedious.
  It is in fact easier to $\eta$-expand both $f$ and $g$ and apply extensionality from \cref{thm:lambda} on the resulting $\lambda$-abstractions.
  In \zflean, this amounts to the following rewriting steps:
  \begin{leanlisting}
  -- hf : IsFunc A B f
  -- hg : IsFunc A B g
  rw [lambda_eta hf, lambda_eta hg, lambda_ext_iff]
  \end{leanlisting}
  This ends up in lifting a proof of extensionality for ZFC functions to standard extensionality of Lean functions.
  Full details of the proof are available in the artifacts.
\end{example}

\subsection{Automation of side-conditions}
Manipulating relations in a set-theoretic style generates repetitive proof obligations regarding relations and partial/total functions.
As already mentioned, the library provides three lightweight tactics whose purpose is to try to solve such obligations \emph{structurally}:
\begin{itemize}
  \item \inlinelean|zrel|: relational goals $R\subseteq A\times B$;
  \item \inlinelean|zpfun|: partial functionality goals $\mathsf{IsPFunc}(f,A,B)$;
  \item \inlinelean|zfun|: total functionality goals $\mathsf{IsFunc}(A,B,f)$.
\end{itemize}
All three tactics work similarly: they first search among local hypotheses for witnesses of the required properties, or apply standard theorems (e.g., composition of functions is a function) to reduce the goal to simpler subgoals.

They may then recursively call each other and backtrack as needed when application of a theorem fails.
In order to reduce search time and keep automation usable at runtime, some theorems are selected and tagged with relevant attributes---called \texttt{zrel}, \texttt{zpfun}, and \texttt{zfun} for consistency with the tactic names---so that they can be found by the corresponding tactic.
This design is illustrated in \cref{fig:tactics}.
This also allows all three tactics to be extended with new rules as needed.
When an automatic discharge fails, it falls back to the user for manual proof.
One can also pass explicit proofs to definitions and theorems instead of relying on automation, even when the tactics would succeed.
%
\begin{figure}[t]
  \centering
  \caption{Overview of the automation tactics for relations and functions. Dashed arrows indicate the shape of the goals discharged by each tactic; thick arrows indicate tactic invocations.}
  \label{fig:tactics}
  \begin{tikzpicture}[
      inner sep=8pt,
      every node/.style={rounded corners=3pt},
      arr/.style={->, shorten <= 3pt, shorten >= 3pt},
    ]
    \small
    \node[draw] (rel_goal) {$\vdash R \subseteq A \times B$};
    \node[draw, left=of rel_goal] (pfun_goal) {$\vdash \mathsf{IsPFunc}(f,A,B)$};
    \node[draw, right=of rel_goal] (fun_goal) {$\vdash \mathsf{IsFunc}(A,B,f)$};
    \node[draw, below=of rel_goal, thick] (zrel) {\inlinelean|zrel|};
    \node[draw, below left=of zrel, thick] (zpfun) {\inlinelean|zpfun|};
    \node[draw, below right=of zrel, thick] (zfun) {\inlinelean|zfun|};

    \draw[arr, dashed] (zrel) -- (rel_goal);
    \draw[arr, dashed] (zpfun) -| (pfun_goal);
    \draw[arr, dashed] (zfun) -| (fun_goal);

    \draw[arr, thick] (zpfun) -- (zrel);
    \draw[arr, thick] (zfun) -- (zrel);
    \draw[arr, thick] (zfun) -- (zpfun);
  \end{tikzpicture}
\end{figure}
%
We show a few examples of typical such theorems (with proofs omitted) in \cref{lst:automation-layer-thm}.
\begin{leanlisting}[title={Examples of theorems tagged for automation.}, label=lst:automation-layer-thm]
@[zrel] -- converse of a relation is a relation
theorem subset_prod_inv {R A B : ZFSet} (hR : R ⊆ A × B) :
  R⁻¹ ⊆ B × A := $\irrelprf$
@[zpfun] -- identity is a partial function
theorem Id.IsPFunc {A : ZFSet} : (𝟙A).IsPFunc A A := $\irrelprf$
@[zfun] -- composition of total functions is a total function
theorem IsFunc_of_composition_IsFunc {g f : ZFSet} {A B C : ZFSet}
  (hg : B.IsFunc C g) (hf : A.IsFunc B f) :
    A.IsFunc C (composition g f A B C) := $\irrelprf$
\end{leanlisting}

Given this infrastructure, we then leverage Lean's \emph{automatic parameters} to pass default tactics in partial definitions depending on side-conditions, so that those side-conditions are automatically discharged when possible.
Automatic parameters are very similar to default parameters, but are specific to tactics: they try to discharge the goal using the provided default tactic if the argument is not supplied explicitly, and fail silently otherwise, leaving the user to provide a manual proof.
Furthermore, those side conditions often depend on parameters that are kept implicit (e.g., the sets $A$ and $B$ in $R \subseteq A \times B$); they are also inferred automatically by Lean.
This ultimately enables working with these definitions---most of the time---without worrying about side conditions or implicit parameters, which is the intended purpose of the \zflean{} framework.

\cref{lst:automation-layer-thm} illustrates this already: the converse relation $R^{-1}$ in theorem \inlinelean|subset_prod_inv| is actually a relation between $B$ and $A$, which are implicit at this level, and relies on the side-condition $R \subseteq A \times B$, automatically discharged by \inlinelean|zrel|.
This principle will be further exemplified in \cref{sec:case-study}.
\begin{leanlisting}[title={Examples of simplification rules for relations and functions.}, label=lst:simp-rules]
@[simp] theorem Image_empty {R A B : ZFSet} (hR : R ⊆ A.prod B) : R[∅] = ∅ := $\irrelprf$
@[simp] theorem range_Id {A : ZFSet} : (𝟙A).Range = A := $\irrelprf$
@[simp] theorem Image_of_composition_inv_self_of_bijective {f A B X : ZFSet}
  {hf : A.IsFunc B f} (hf : f.IsBijective) (hX : X ⊆ A) : f⁻¹[f[X]] = X := $\irrelprf$
\end{leanlisting}
We also tag relevant equality theorems with the \texttt{simp} attribute, so that they can be used by Lean's \inlinelean|simp| tactic for rewriting.
A few examples are shown in \cref{lst:simp-rules}, with proofs omitted.

\iffalse
  \begin{table}[t]
    \centering
    \caption{Summary of the relational interface provided by the library and used in the paper. Implicit parameters are omitted, and automatic parameters are shown separately.}
    \label{tab:rel-api}
    \small
    \begin{tabular}{lp{.33\columnwidth}p{.33\columnwidth}}
      \hline
      \textbf{Symbol}                                               & \textbf{Description} & \textbf{Auto. param.} \\
      \hline
      $\Id{A}$                                                      &
      Identity relation on $A$.                                     &
      \\

      $R^{-1}$                                                      &
      Converse of a relation.                                       &
      $R \subseteq A \times B$                                                                                     \\

      $\Dom(R)$                                                     &
      Domain of a relation.                                         &
      $R \subseteq A \times B$                                                                                     \\

      $\Ran(R)$                                                     &
      Range of a relation.                                          &
      $R \subseteq A \times B$                                                                                     \\

      $\Image{R}{X}$                                                &
      Image of a set under a relation.                              &
      $R \subseteq A \times B$                                                                                     \\

      $\fapply{f}{\langle x, h_x \rangle}$                          &
      Function application.                                         &
      $\mathsf{IsPFunc}(f, A, B)$                                                                                  \\

      $g \zbinop{\circ} f$                                          &
      Composition of total functions.                               &
      $\mathsf{IsFunc}(A, B, f)$,
      $\mathsf{IsFunc}(B, C, g)$                                                                                   \\

      $\zlambda{\!A\!}{\!B}{\!x\!}{\!e(x)}$                         &
      $\lambda$-abstraction.                                        &
      \\

      $A \zbinop{\hookrightarrow} B$, $B \zbinop{\hookleftarrow} A$ &
      $A$ can be embedded into $B$.                                 &
      \\

      $A \zbinop{\cong} B$                                          &
      $A$ is isomorphic to $B$.                                     &
      \\
      \hline
    \end{tabular}
  \end{table}
\fi

Overall, this calculus provides a compact, rewriting-friendly and easily extensible interface for extensional reasoning with relations and functions \emph{inside} ZFC, while hiding routine side-conditions behind small, domain-specific tactics.
% \cref{tab:rel-api} summarizes the relational interface provided by the library, notations and the associated automatic parameters.
% \fixme[Does the table really add value to the paper?]

